\chapter{Rock: Turning Current into a Heartbeat}
\label{ch:genre-rock-deep}

\classquote{审一以定和，比物以饰节，节奏合以成文。}{《礼记·乐记》}{Music first fixes one tone to establish harmony, then uses different instruments to ornament the beat; rhythm joined together becomes the pattern.}

\section{What rock is: a skeleton first, then current}
\label{sec:rock-what}

\term{rock}{rock} is not a synonym for ``turn the volume up''. It is a way of pushing body and emotion forward at the same time: a riff you can remember, a clear backbeat, a layer of weighty distortion, and a release that actually arrives. It can be rough or refined; it can be three people or a band with keyboards, brass and a second guitar. What makes it work is that the skeleton stands first, so the current has somewhere to travel.

Four handles:

\begin{longtable}{p{0.16\textwidth}p{0.32\textwidth}p{0.38\textwidth}}
\toprule
Handle & What it feels like & The writing question \\
\midrule
\endfirsthead
\toprule
Handle & What it feels like & The writing question \\
\midrule
\endhead
Riff & A short figure that repeats like the song's ID card & If you remove the riff, does the song still recognise itself? \\
Backbeat & Snare on 2 and 4; the body is pushed along & Is the backbeat clear? Is the kick solid on 1 and 3? \\
Distortion & Guitar/bass harmonics thicken; emotion gains a ``current'' & Is the distortion carrying the main riff, or just blurring every track together? \\
Release & The verse holds back; the chorus or section end truly opens & What does the chorus add? If it is only louder, is that enough? \\
\bottomrule
\end{longtable}

Rock is often most moving not when it is loudest, but when it has held back for a long time and finally lets go. A quiet verse, a bridge with only bass and drums, two beats of silence before a solo: all of these exist so the current has contrast. Without holding back, there is no release.

\section{From garage to stadium: how rock got here}
\label{sec:rock-history}

Rock's lineage is not a straight line. It grew out of blues, country, rhythm and blues and gospel, then each generation rewrote it with its own equipment and venues. The table follows one thread only: when instrumentation and media change, the writing changes with them.

\begin{longtable}{p{0.16\textwidth}p{0.34\textwidth}p{0.42\textwidth}}
\toprule
Period & Key shift & Effect on arrangement \\
\midrule
\endfirsthead
\toprule
Period & Key shift & Effect on arrangement \\
\midrule
\endhead
1950s & Rock and roll enters the teenage market; Chuck Berry, Elvis Presley and Little Richard push electric guitar, piano and backbeat to the front & Beats 2 and 4 become a body coordinate; electric guitar rises from accompaniment to signature; songs are short and hooks direct \\
1960s & The British Invasion (The Beatles, The Rolling Stones), Jimi Hendrix's feedback and distortion, The Who's high-energy shows & Guitar timbre becomes compositional material; the studio turns from documenting a performance into manufacturing sound; albums hold longer structures \\
1970s & Led Zeppelin and Black Sabbath's heavy riffs; Pink Floyd's progressive structures; Ramones and Sex Pistols' punk backlash & Heavy distortion, slow/medium riffs, long sections and three-minute punk coexist; stage volume and amplifiers become part of the instrument \\
1980s & MTV, metal, alternative rock, U2 and R.E.M.; drum machines and synths enter rock & Image and timbre shape band identity together; guitar walls and large-room reverb become ways to sound ``big'' \\
1990s & Nirvana, Pearl Jam and Soundgarden's grunge; Radiohead and other alternative rock & ``Quiet--loud'' dynamics become structure; rough timbres and real drums return; lyric and attitude matter as much as sound \\
2000s & The White Stripes, Arctic Monkeys, garage/indie rock revival; digital release & Small line-ups, strong riffs, retro recording texture; two or three people can still make a complete low end and impact \\
2010s--present & Streaming, festivals, genre mixing; rock borrows from electronic, hip-hop and pop & Riffs can be short, timbres can be ``non-guitar''; form is more fragmented, but live release remains the core \\
\bottomrule
\end{longtable}

Chinese rock is often traced to the 1980s: Cui Jian's ``Nothing to My Name'' in 1986 brought electric guitar, a band and individual expression into the public view; in the 1990s Tang Dynasty, Black Panther and others turned heavy metal and hard rock into an industry; from the 2000s, indie bands, festivals and the internet pushed rock into more fragmented scenes. Industrial conditions differ by region, but the question is the same: how does a limited line-up produce a clear backbeat, riff and release?

\section{Arrangement: every instrument bites the same skeleton}
\label{sec:rock-craft}

Rock can be sparse, but it cannot be ``everyone playing their own song''. Writing roles into track names is more useful than filling the project with amp simulators.

\subsection{Typical instrumentation and five layers}
\label{sec:rock-instrumentation}

\begin{longtable}{p{0.16\textwidth}p{0.25\textwidth}p{0.25\textwidth}p{0.20\textwidth}}
\toprule
Layer & Common instruments/tracks & Main job & Common mistake \\
\midrule
\endfirsthead
\toprule
Layer & Common instruments/tracks & Main job & Common mistake \\
\midrule
\endhead
Voice & Lead, harmony, shout & Tell the story, carry attitude, own the chorus memory & Too much tuning/reverb turns a rock voice into pop \\
Rhythm guitar & Distorted/semi-distorted guitar, power chords & Provide the riff, the chord wall, lock the low end with bass & Wall-to-wall strumming leaves the voice no room to breathe \\
Lead guitar & Distorted guitar, slide/feedback & Short answers, solos, section-end hooks & Solos chase speed without motif or space \\
Bass & Electric bass, overdrive/distortion & Lock the kick, fill roots, give the low end a bite & Doubling the guitar exactly turns the low end into a wall \\
Drums & Acoustic kit, percussion & Set the backbeat, set the section floor, provide fills & Filling every section leaves the chorus with no arrival \\
\bottomrule
\end{longtable}

A practical starting point is six tracks: lead vocal, rhythm guitar, lead guitar, bass, drums, and the master bus. Run the five-instrument relationship first, then add a second guitar, keys or percussion. Track count is not a rock medal; every instrument knowing which beat it bites is.

\subsection{Groove: clear backbeat, solid kick}
\label{sec:rock-groove}

The most common rock skeleton is 4/4 at 90--140 BPM; punk is faster, heavy/doom is slower. Whatever the tempo, the snare on 2 and 4 must be clear, and the kick usually lands on 1 and 3 or together with the bass's key riff.

Three rules:

\begin{itemize}
\item \textbf{Lock kick and bass first}: the bass root lands on the kick, and the riff's strong note can land with both; do not make the bass chase the guitar on every beat.
\item \textbf{One owner for the backbeat}: the snare already owns 2 and 4; if the rhythm guitar accents every 2 and 4 in the same way, the backbeat blurs.
\item \textbf{Save fills for section ends}: put drum fills at verse-to-chorus, solo-to-chorus or the outro; a fill every four bars exhausts the listener.
\end{itemize}

\subsection{Guitar: one riff, three jobs}
\label{sec:rock-guitar}

Electric guitar usually does three jobs in rock; do not make it do all three at once:

\begin{enumerate}
\item \textbf{Riff}: short, low, repeating, usually a \term{power chord}{power-chord} (root plus fifth) or a single note; it owns the memory point.
\item \textbf{Comping}: give chords at the end of vocal phrases, or strum eighths in the intro/chorus; it owns harmony and push.
\item \textbf{Lead}: short answers, solos, feedback/noise; it owns ``speaking'' and release.
\end{enumerate}

More distortion is not automatically better. Give the main riff one distortion setting; the rhythm guitar can use a slightly cleaner tone for contrast; the lead guitar steps back in the chorus. Two distorted guitars stacked from top to tail blur the low end and remove the voice's place.

\subsection{Bass and drums: the low-end bite}
\label{sec:rock-rhythm-section}

Bass is not ``guitar an octave lower''. It builds the low-end floor with the kick:

\begin{itemize}
\item roots land on strong beats; passing notes only in the beat before a chord change;
\item if the guitar riff sits high, the bass can carry a low-register counterpart instead of doubling it exactly;
\item distorted bass can add midrange bite, but do not push the low end until it cannot be separated from the kick;
\item the kick should be short and solid; the snare needs a backbeat ``snap''; do not let reverb soften the backbeat.
\end{itemize}

\subsection{Voice and form: hold, then release}
\label{sec:rock-vocal-form}

A rock voice does not have to shout all the time. The verse can use a lower chest voice, close miking and little reverb; the chorus opens the force and the space. The key is to make the opening have contrast.

Common forms:

\begin{itemize}
\item Intro riff → verse (riff + voice) → chorus (thicker guitar, heavier drums) → verse → chorus → bridge (new harmony/lower density) → solo → chorus → outro.
\item The riff can also mark sections: intro, chorus entrance and outro all use the same riff, so the listener knows ``it is back''.
\item A solo is not ``more notes''; it amplifies the verse's motif. The solo section can keep only lead guitar, bass and drums, giving the ear room.
\end{itemize}

\section{Four songs, arrangement only}
\label{sec:rock-cases}

The analyses below study arrangement only. They quote no lyrics and include no melodic notation. Work information follows album release information and official video pages; links are in the chapter's further reading.

\subsection{Queen, ``Bohemian Rhapsody'': a song written as several acts}
\label{sec:rock-case-bohemian}

Queen's ``Bohemian Rhapsody'' was released in 1975 on \emph{A Night at the Opera}. It does not carry the whole song on one riff; it assembles piano, choir, guitar and operatic sections into a ``multi-act'' piece.

\begin{itemize}
\item \textbf{Form}: piano/choir intro, narrative section, guitar section, opera section, hard-rock section, outro; each has a different density and timbre, so the listener is led rather than waiting for one chorus.
\item \textbf{Orchestration}: piano carries the harmonic skeleton, stacked voices create the ``choir'', and the guitar only truly pushes in the hard-rock section; each instrument is the lead in its own act.
\item \textbf{Logic}: section contrast replaces loop stacking. This differs from standard pop/rock writing, but it proves one principle: a riff is not the only memory point; a section can be one too.
\end{itemize}

Takeaway: if your song has several different emotional states, do not force them into one repeating chorus. Design an anchor---a main riff or phrase to return to---so the listener has a way back through the change.

\subsection{Nirvana, ``Smells Like Teen Spirit'': quiet--loud is the structure}
\label{sec:rock-case-teen-spirit}

Nirvana's ``Smells Like Teen Spirit'' was released in 1991 on \emph{Nevermind}. What is most useful to learn is the dynamic design between a quiet verse and an explosive chorus.

\begin{itemize}
\item \textbf{Guitar}: the main riff uses a low register, power chords and distortion; the verse holds it down, the chorus amplifies the same motif instead of replacing it with new material.
\item \textbf{Drums}: snare on 2 and 4 and kick create a clear backbeat; crash and denser drums in the chorus give the release physical weight.
\item \textbf{Voice}: the verse is close, low and dry; the chorus opens without over-tuning; strong lyric words land on backbeats/strong beats.
\item \textbf{Logic}: the same material changes density, distortion and vocal force to make ``quiet--loud''; the chorus does not need many extra tracks to explode.
\end{itemize}

Takeaway: write a riff that works at low density first, then decide what the chorus adds. If the verse has already used everything, the chorus can only be louder, not more ``open''.

\subsection{The White Stripes, ``Seven Nation Army'': one riff, two people, one stadium}
\label{sec:rock-case-seven-nation}

The White Stripes' ``Seven Nation Army'' was released in 2003 on \emph{Elephant}. It uses one low-register riff, minimal drums and two people to make a stadium-sized memory point.

\begin{itemize}
\item \textbf{Riff}: the bass-like guitar riff is the main hook; it is low, short and repeating, and tens of thousands of people can sing it live.
\item \textbf{Line-up}: only guitar, drums and voice, with no bass; the riff carries the low end itself, and the drums pin the backbeat with a minimal kick/snare.
\item \textbf{Space}: the verse leaves room; the chorus adds vocal/guitar layers; it does not add more instruments but repeats the same riff at different densities.
\item \textbf{Logic}: a small line-up is not a limitation; it enlarges each sound's responsibility. Clear riff, clear drums, clear voice, and the song stands.
\end{itemize}

Takeaway: if you have only two or three tracks, do not rush to fill them with a pad. Ask first whether the riff can carry both memory and low end. If it can, make it clear.

\subsection{Arctic Monkeys, ``Do I Wanna Know?'': a riff can be cold}
\label{sec:rock-case-do-i-wanna-know}

Arctic Monkeys' ``Do I Wanna Know?'' was released in 2013 on \emph{AM}. It makes the rock riff cold, slow and nocturnal rather than relying on high-gain distortion.

\begin{itemize}
\item \textbf{Riff}: a low-register guitar motif repeats with a little slide and open strings; its ``cool'' comes from timbre and space, not speed.
\item \textbf{Groove}: the backbeat is clear but unhurried, and the bass follows the riff's low notes; the whole feels like a half-time sway, leaving room for the voice.
\item \textbf{Voice}: the lead is close and low, the chorus opens with falsetto/harmony; the voice does not shout, it holds the emotion inside.
\item \textbf{Logic}: the riff owns memory, drums/bass own sway, voice owns emotion; each does one job, so it sounds empty without being thin.
\end{itemize}

Takeaway: distortion is not rock's only answer. A narrow riff timbre, a clear low end and a solid backbeat can be just as rock.

\section{The yeban skeleton: eight bars of rock}
\label{sec:rock-yeban}

The skeleton below can be rebuilt in yeban. The focus is riff, backbeat and ``quiet--loud''.

\begin{longtable}{p{0.10\textwidth}p{0.22\textwidth}p{0.58\textwidth}}
\toprule
Bar & Section & Arranging action \\
\midrule
\endfirsthead
\toprule
Bar & Section & Arranging action \\
\midrule
\endhead
1--4 & Verse & Rhythm guitar: low-register riff, power chord/single note; bass: root locks the kick; drums: kick 1/3, snare 2/4, closed hats in eighths; voice: middle-low register, space at phrase ends; lead guitar: out \\
5--8 & Chorus & Rhythm guitar: riff plus eighth strumming/double track; bass: root plus passing note; drums: crash and denser fill; voice: open, add harmony; lead guitar: short answers, never cover the voice \\
\bottomrule
\end{longtable}

Order of operations:

\begin{enumerate}
\item Create a 100 BPM, 4/4 project; build lead vocal, rhythm guitar, lead guitar, bass and drums, and keep the master as a bus.
\item Write the rhythm guitar's two-bar riff (for example, an E5--G5--A5 power-chord loop), keeping it low; use F7's piano roll to fix each beat's landing.
\item Write drums: kick 1/3, snare 2/4, closed hats in eighths; add a two-beat fill at the end of bar 4 into the chorus.
\item Write bass: roots follow the kick; the riff's strong notes land with the guitar; do not chase the guitar on every beat.
\item Write the lead: verse stays middle-low, with half a beat to a beat of space at phrase ends; the chorus opens on bar 5, but do not push every line to the top.
\item On bar 5, add the lead guitar's short answer, crash and one harmony layer; add only one layer and keep the lead clear.
\item Use F8 to check levels; solo each track and inspect the low end: bass and kick must bite, rhythm and lead guitar must separate in the midrange.
\item Use F3 for the stacked view to compare timbres; if distortion covers the voice, reduce distortion/rhythm guitar first instead of only raising the vocal.
\item Save the project and export an eight-bar audition; ask ``did the verse hold back? did the chorus really open?''
\end{enumerate}

\section{Exercises: three paths, different trade-offs}
\label{sec:rock-practice}

\begin{exercise}
\textbf{A two-bar riff, three versions.} Using the same E--G--A, make three versions: A a low single-note riff; B power chords (root plus fifth); C power chords with palm muting. Compare ``memory point, low-end weight, room left for the voice''. A is clear but may be thin; B has weight but blurs easily; C bites but needs control of mute length.
\end{exercise}

\begin{exercise}
\textbf{Quiet--loud, three lifts.} Attach three choruses to the same verse: A only fill out the drums; B only open the guitar distortion; C remove the lead guitar/harmony first, then bring everything back on the chorus downbeat. Play each and note which ``opening'' has the most contrast; try all three before choosing.
\end{exercise}

\begin{exercise}
\textbf{The riff and the drums bite.} Put the same riff against the kick on 1/3, against the backbeat on 2/4, and against a weak sixteenth-note position. Change only the drums, never the riff. Hear which feels most ``push'' and which most ``drag''. Write the result in one sentence in the project notes.
\end{exercise}

\section{Further reading}
\label{sec:rock-reading}

\begin{sloppypar}
\begin{itemize}
\item \textbf{AllMusic, ``Pop/Rock''}: \url{https://www.allmusic.com/genre/pop-rock-ma0000002613}. Use the genre entry to fill in rock's lineage, subgenres and representative artists; good for building a coordinate system first.
\item \textbf{Queen, ``Bohemian Rhapsody'' (Official Video Remastered)}: \url{https://www.youtube.com/watch?v=fJ9rUzIMcZQ}. See how one song uses multiple sections and timbres to create a ``multi-act'' arrangement.
\item \textbf{Nirvana, ``Smells Like Teen Spirit'' (Official Music Video)}: \url{https://www.youtube.com/watch?v=hTWKbfoikeg}. Compare density, distortion and drum layers between the quiet verse and the explosive chorus.
\item \textbf{The White Stripes, ``Seven Nation Army'' (Official Music Video)}: \url{https://www.youtube.com/watch?v=0J2QdDbelmY}. Hear how two people use one riff and minimal drums to carry a complete low end and memory point.
\end{itemize}
\end{sloppypar}

The links above were checked one by one against official pages and search results for title and author on 27 November 2026. If a platform's regional restriction prevents opening one, search for the entry name on the artist's official account in that platform.

\section{Chapter takeaways}
\label{sec:rock-takeaways}

\begin{itemize}
\item Rock is not volume; it is the combination of riff, backbeat, distortion and release.
\item Rock's evolution follows media and live performance: radio, TV, MTV, streaming and festivals each changed what sound can represent a band.
\item Write function before instrument; rhythm guitar, lead guitar, bass and drums each do one job.
\item The backbeat must be clear and the kick solid; save fills for section ends, not every four bars.
\item More distortion is not automatically better; give the main riff weight and leave the voice a place.
\item Quiet--loud is one of rock's most reliable structures: hold the verse down so the chorus has something to release.
\item The yeban skeleton runs ``riff---drums/bass---voice---lead/colour---level check'' in that order; editable and explainable is good arrangement.
\end{itemize}

\begin{helpmapping}
\begin{longtable}{p{0.45\textwidth}p{0.45\textwidth}}
\toprule
Section & Help anchor \\
\midrule
What rock is & \texttt{arrangement-genre-rock-what} \\
From garage to stadium & \texttt{arrangement-genre-rock-history} \\
Arrangement: every instrument bites the same skeleton & \texttt{arrangement-genre-rock-craft} \\
Four songs, arrangement only & \texttt{arrangement-genre-rock-cases} \\
The yeban skeleton & \texttt{arrangement-genre-rock-yeban} \\
Exercises & \texttt{arrangement-genre-rock-practice} \\
Further reading & \texttt{arrangement-genre-rock-reading} \\
\bottomrule
\end{longtable}
\end{helpmapping}
